\section{Introduction}\label{sec:intro}

An attractor is called hidden if its basin of attraction does not intersect small neighborhoods of equilibria; otherwise, the attractor is self-excited.\cite{LeonovKuznetsov2013,Leonov2011PLA,Kuznetsov2020JCSSI} Hidden oscillations occur in Chua circuits,\cite{Leonov2012PhysD,Stankevich2017,Kuznetsov2023Chua} in control systems with saturation,\cite{Andrievsky2017Aircraft} in the Aizerman and Kalman problems,\cite{Bragin2011} in generalized Lorenz and Rabinovich systems,\cite{Chen2017OnePath,Kuznetsov2018Rabinovich,Kuznetsov2020Lorenz} in systems without equilibria,\cite{JafariSprottGolpayegani2013,Wei2011} with a single stable equilibrium,\cite{WangChen2012,Molaie2013} with lines and surfaces of equilibria,\cite{JafariSprott2013,Jafari2016Surfaces} in two-dimensional maps,\cite{Jiang2016Maps,WangDing2018Map,Huynh2019} and in neural network models.\cite{DancaKuznetsov2017Hopfield} Reviews\cite{Dudkowski2016,JafariSprottNazarimehr2015,GuanXie2025} show that the number of such examples has grown by orders of magnitude over the past fifteen years. Since a hidden attractor usually coexists with other attractors, searching for it is a problem of multistable dynamics.\cite{Feudel2008,PisarchikFeudel2014} Attractors whose basins occupy a small fraction of the state space are called rare.\cite{Chudzik2011Rare,Yue2021RareDuffing}

The main practical difficulty stems from the definition itself. The standard procedure for finding attractors launches trajectories from neighborhoods of unstable equilibria along their unstable manifolds. By definition, this procedure does not find a hidden attractor, and with an unlucky choice of initial conditions, such an attractor is easily missed. Analytical--numerical methods have been proposed for localizing hidden attractors: harmonic balance and the describing function followed by continuation in a parameter.\cite{Leonov2011PLA,KuznetsovLeonovVagaitsev2010,Kuznetsov2017IFAC} There are also geometric heuristics: perpetual points,\cite{Prasad2015,Dudkowski2015PP,Dudkowski2017PP} connecting curves\cite{Gilmore2010CC,GuanXie2021CC,Naseri2021} and surfaces,\cite{Song2023CS} coordinate-wise zero-velocity sets,\cite{Godara2018CVS} and parameter switching.\cite{Danca2018Switching} Evolutionary search over initial conditions is used as well.\cite{Zelinka2016,Zelinka2021Chapter} A general tool is the estimation of basins by sampling initial conditions,\cite{Menck2013,DatserisWagemakers2022,Datseris2023Framework,StenderHoffmann2022,Gelbrecht2020MCBB} together with set-oriented methods.\cite{Hsu1987,DellnitzHohmann1997} A sample of $n$ initial conditions detects an attractor of basin fraction $p$ with probability $1-(1-p)^n$,\cite{Datseris2023Framework} so rare attractors are found either by very large samples or by continuation from parameter values where they are not rare: path following combined with basin stability\cite{Brzeski2017} and global continuation, which seeds the search at each parameter value with the attractors found at the previous one.\cite{Datseris2023Framework,Datseris2026Global,Haerter2026BBT} The limitations of the heuristics are well known. Perpetual points may not exist at all, or they may lead to something other than the target.\cite{Nazarimehr2017PP,Jafari2015LimitPP,Jafari2026NoPP} Connecting curves do not always pass through the basin of the hidden attractor.\cite{GuanXie2025}

Assessments of localization methods almost always remain qualitative: ``the attractor was found from such-and-such a point.'' How many integrations were needed, with what probability the method finds the target within a given budget, and how much better it is than uniform sampling are usually not reported. This picture emerges from the review in Ref.~\onlinecite{GuanXie2025} and from an analysis of the full text of Ref.~\onlinecite{Zelinka2016}. Meanwhile, hiddenness and rarity are different properties. In well-known examples, the basin of a hidden attractor often occupies most of the natural region of initial conditions; for the Rabinovich system at $r=100$, it occupies 76\% (Sec.~\ref{sec:bench}). What makes the problem hard is a small basin measure $p$, not the absence of a connection to equilibria. In addition, a long-lived transient chaotic set is sometimes declared a ``hidden attractor.'' Examples are the story of the ``hidden attractor'' of the classical Lorenz system\cite{Munmuangsaen2018,SprottMunmuangsaen2018Comment} and the Rabinovich--Fabrikant system;\cite{Danca2016Transient} see also Refs.~\onlinecite{YorkeYorke1979,GrebogiOttYorke1983,LaiTel2011}. Comparing methods therefore requires an explicit verification protocol.

If the target basin is small, finding the attractor becomes the problem of detecting a rare event in the space of initial conditions. Well-developed tools exist for such problems. They include multilevel splitting\cite{Glasserman1999,CerouGuyader2007,Cerou2019AMS} and its static counterpart, subset simulation,\cite{AuBeck2001} the cross-entropy method,\cite{Rubinstein1999CE} and sequential design of experiments with active learning of the failure domain.\cite{Bect2012,Echard2011AKMCS} Novelty search and quality-diversity algorithms belong here as well.\cite{LehmanStanley2011,MouretClune2015,Pugh2016QD,CullyDemiris2018} Rare-event algorithms have been applied, for example, to extreme events and transitions in climate models\cite{Ragone2018,Lestang2018TAMS} and to high-dimensional dynamical systems.\cite{Baars2021AMS} Machine-learning methods for basins of attraction\cite{Wang2022HAL,Cao2025AL,Yang2024PRR,DNNBasins2021,Gameiro2025,Haerter2026BBT} estimate the fractions of already known basins rather than search for unknown attractors. To the best of our knowledge, these tools have not been applied systematically to the search for hidden attractors.

The contributions of this work are as follows.
\begin{enumerate}
\item The search for a hidden attractor is formulated as a rare-event detection problem. The quality of a method is measured by the detection probability $P(N)$ as a function of the budget $N$, by the cost $N_{0.9}$, by the gain factor $G$ relative to uniform sampling, and by the scaling exponent $s$ in $N_{0.9}\propto p^{-s}$. We identify three channels of information through which an adaptive method can outperform uniform sampling at all.
\item We build a four-level benchmark: three-dimensional flows, maps, four-dimensional systems, and synthetic families with analytically known basin measure. We develop a protocol for identifying and verifying attractors. While refining it, we found a numerical artifact in the Chua circuit and a discrepancy with a published table of basin fractions.
\item We propose a portfolio of adaptive strategies: multilevel splitting driven by transient duration, active learning of basin boundaries, and a uniform share. We prove that in the worst case it is at most $1/\varepsilon$ times costlier than uniform sampling and inherits the scaling exponent of the splitting arm. We construct a cost model with diagnostics that predict when a gain is possible.
\item We measure the gain and its scaling with basin measure and dimension, and we measure the cost not only in integrations but in integrator steps and wall time, against uniform sampling with the same evaluation and against uniform sampling with recurrence-based early stopping (Attractors.jl). The gain carries over to time where trajectories are bounded and shrinks where most of them escape quickly. On two coupled Chua circuits, a rare hidden regime ($p=2\times10^{-5}$) is found 55 times faster in integrations and 10--64 times faster in time, and the enumeration of regimes by the identification protocol reveals a hidden chaotic attractor of the coupled circuits that is born by the coupling and that we have not found reported. Failure cases, including the limit set by the time resolution of the reaction coordinate, are documented explicitly; enlarged regions of initial conditions and continuation in a parameter do not make the hidden attractors of the benchmark rare.
\end{enumerate}
